% demo-mysub.tex -- mysub / mysub* environments and \mysubfigure
\documentclass[sectioning=nobreak]{myreport}
\usepackage{mydemo}
\begin{document}
\demoheader{mysub, \string\mysubfigure}{A figure made of several
sub-figures, each with an optional sub-caption (a), (b), ...}

\mysection*{Syntax}
\begin{demosyntax}
\begin{mysub}{<label>}[<short caption>]{<caption>}[<position>]
  \mysubfigure[<options>]{<file>}[<sub-caption>]{<width>}[<alignment>]
  ...
\end{mysub}
\begin{mysub*}{<label>}{<caption>}[<position>]   % not in the list of figures
\end{demosyntax}
\begin{itemize}
  \item \texttt{<label>}: the figure is \texttt{fig:<label>}, its
    sub-figures \texttt{fig:<label>:a}, \texttt{fig:<label>:b}, \dots
  \item \texttt{<width>} of \verb|\mysubfigure|: fraction of
    \verb|\textwidth| given to the sub-figure; \texttt{<options>}
    (\texttt{width=}, \texttt{height=}) are relative to that width
    (default \texttt{width=1}).
  \item \texttt{<alignment>}: vertical alignment of the sub-figure
    (\texttt{T} by default, or \texttt{t}, \texttt{c}, \texttt{b},
    \texttt{B}).
\end{itemize}

\mysection*{Examples}
\begin{demo}
\begin{mysub}{pair}{Two sub-figures side by side.}[H]
  \mysubfigure{wide.png}[A wide image]{0.48}
  \hfill
  \mysubfigure[height=0.1]{tall.png}[A tall image]{0.48}
\end{mysub}
See \myref{fig:pair} and \myref{fig:pair:b}.
\end{demo}

\begin{demo}
\begin{mysub}{three}[Three images]{Three sub-figures, the last one
without sub-caption, bottom-aligned.}[H]
  \mysubfigure{wide.png}[A]{0.3}[b]\hfill
  \mysubfigure{square.jpg}[B]{0.3}[b]\hfill
  \mysubfigure[width=0.5]{tall.png}{0.3}[b]
\end{mysub}
\end{demo}

\begin{demo}
\begin{mysub*}{hidden}{A starred group: not in the list of figures.}[H]
  \mysubfigure{square.jpg}[Only one]{0.3}
\end{mysub*}
\end{demo}
\end{document}
